\chapter*{Abstract}
\label{abstract}

% SOURCE: results_abstract_v1.md
% STATUS: converted from markdown
% NOTE: [source:] tags stripped per ASSEMBLY_PLAN.md (harvested to appF)


Markerless injury-risk screening systems require validated measurement accuracy before deployment in clinical environments, but are typically implemented on internal-consistency evidence alone. This dissertation presents a monocular, markerless posture tracking pipeline evaluated across three sagittal athletic movements (squats, lunges, and drop-jumps) utilizing consumer-grade 30-Hz video. Synchronized 3D optoelectronic and force-plate references from a dynamic drop-jump landing cohort (8 subjects, 48 trials) serve as the physical validation benchmark. Ground-truth comparison reveals a constant peak deep-flexion overestimation bias of +10.52° (95\% Limits of Agreement [LoA]: [-5.54°, 26.58°]) that is depth-independent (Pearson r = -0.16, p = 0.13), contrasting with a shallow contact underestimation bias of -6.69° (95\% LoA: [-26.77°, 13.39°]). For movement screening, incorrect repetitions are characterized by deeper peak flexion depth in squats (Cohen's d = 1.7306, n = 98 reps) and lunges (d = 1.6904, n = 61 reps). Cross-exercise analysis demonstrates propulsive ascent-velocity divergence in lunges (peak ascent d = -0.9721) that is absent in squats. Leave-one-subject-out sequence modeling confirms endpoint biomarker dominance, as deep LSTMs collapse to majority-class guessing (50.00\% balanced accuracy) when amplitude is stripped, whereas a single peak flexion baseline achieves 81.36\% balanced accuracy on squats. The uncertainty framework transfers validation-derived projection variance to slower sagittal tasks, yielding a peak flexion transfer weight of 57.15\%.

Scoped strictly to kinematic screening rather than clinical injury prediction, this work evaluates small analytical cohorts (9 squat subjects, 7 lunge subjects) with architectural demonstrations of personalized noise floors. This dissertation establishes a defensible screening spine through four contributions: cross-exercise integration, a categorized failure-mode taxonomy, uncertainty-weighted transfer of validation errors, and exact-margin counterfactual explanations faithful by construction. This system provides a mathematically bounded foundation for scaling non-invasive clinical movement screening.